\section{Environmental impacts}
\label{sec:environmental_impacts}

The environmental impacts of \acp{bss} have been a central part of their justification~\cite{Shaheen2010, Fishman2016, Midgley2011}. These impacts extend beyond avoided transport emissions and include energy use, material consumption, logistics, maintenance, and end-of-life waste and recycling. While early academic discussions often focused on \ac{ghg} reductions from substituting car trips, recent research has broadened the perspective to include the full life cycle of shared bicycles, docking infrastructure, batteries, and logistics operations.

    \subsection{Avoided transport emissions through modal substitution} 

    Transport emissions comprise both \acp{ghg} (e.g., \ce{CO2}, \ce{CH4}, \ce{N2O}), which accumulate in the atmosphere and drive climate change, and local air pollutants (e.g., \ce{NOx}, \ce{HC}, \ce{PM_{2.5}}, \ce{PM_{10}}, \ce{CO}), which are short-lived and affect urban air quality. Early studies typically estimated avoided emissions by combining (i) cycling distance, (ii) a counterfactual mode that would otherwise have been used, and (iii) mode-specific emission factors, with cycling treated as negligible at the tailpipe. When this approach was combined with simple distance rules for modal substitution and strong extrapolation to the city or fleet scale, it produced very large headline \ce{CO2} savings in systems such as Barcelona's \textit{Bicing}~\cite{Rojas-Rueda2011}, Shanghai's \textit{Mobike}~\cite{Zhang2018}, Beijing's public system~\cite{Qiu2018}, and New York's \textit{Citi Bike}~\cite{Chen2022c}. The savings magnitude is sensitive to how distance is measured (e.g., straight-line vs.\ network routing), but it is mainly driven by substitution assumptions: multi-city surveys typically find car replacement on the order of 5--20\% of trips~\cite{Murphy2015, Fishman2013, Fishman2014}, whereas some published scenarios assumed much higher shares, including 75\%~\cite{Qiu2018} and 90\%~\cite{Rojas-Rueda2011}. If walking or \ac{pt} would have been the true counterfactual, assigning a car counterfactual instead inflates avoided \ce{CO2}, because the emissions difference between cycling and walking or \ac{pt} is much smaller than between cycling and private motorized travel.

    Accordingly, subsequent studies have sought to ground substitution in observed behavior and travel surveys rather than fixed rules, in line with the evidence reviewed in Section~\ref{sec:mobility_and_transport_impacts}. Kou et al.~\cite{Kou2020} developed a trip-level probabilistic assignment that combines observed bike-sharing distances with city-specific travel surveys and life-cycle emission factors, assigning each trip to a most likely replaced mode; applied to eight U.S. cities in 2016, it yields far smaller annual savings than many previous assessments. Saltykova~\cite{Saltykova2022} used traveler surveys for a dockless system in Chengdu, China, finding that 39\% of trips replaced bus, 13.5\% replaced subway, and only 10.9\% replaced car trips. With these figures, avoided emissions fell markedly compared to a ``cars-only'' counterfactual, illustrating how omitting \ac{pt} can inflate climate benefits by a factor of two or more. 
    
    A more refined approach was proposed by Li et al.~\cite{Liu2021}, who developed a high-resolution discrete-choice framework matching over 23 million dockless bike-sharing trips in Shanghai to realistic modal alternatives using travel time, distance, and cost information from online routing services. Their results showed that most trips substituted for walking (30.7\%) or \ac{pt} (50\%), with car replacement below 2\%. This progression from simple distance rules to survey-informed probabilities, and finally to discrete-choice models, illustrates a gradual shift toward more behaviorally grounded and context-sensitive estimates of environmental benefits.
    
    Research has also estimated reductions in local air pollutants (e.g., \ce{NOx}, \ce{SO2}, particulate matter) following the same logic as with \ac{ghg} emissions, though such studies remain relatively uncommon and inherit the same substitution caveats as \acp{ghg} accounts~\cite{Zhang2018, Qiu2018}.



    \subsection[LCA for BSSs]{\ac{lca} for \acp{bss}}

    Avoided transport emissions capture only part of the environmental picture, and \acp{bss} operations can even offset expected benefits: Fishman, Washington, and Haworth~\cite{Fishman2014} showed that London’s \ac{bss} in 2012 increased overall motor vehicle usage due to rebalancing. Consequently, recent work relies on \acp{lca}--typically aligned with ISO 14,040 and 14,044~\cite{LCA2006a, LCA2006b}--to account for burdens from manufacturing, operation, maintenance, and end-of-life management, while also incorporating modal substitution within a single framework (Figure~\ref{fig:LCA_bss}). Results are usually reported per service delivered as g~\ce{CO2}-eq/pkm~\cite{Luo2019, Chen2024, Bonilla-Alicea2020}, where \ce{CO2}-eq aggregates the warming effect of multiple \acp{ghg} (e.g., \ce{CO2}, \ce{CH4}, \ce{N2O}) into an equivalent amount of \ce{CO2} based on global warming potentials. Reported values for complete \ac{bss} \acp{lca} span roughly 30--150~g~\ce{CO2}-eq/pkm~\cite{Zhu2023, Felipe-Falgas2022, Cazzola2020, Sun2022, Calan2024}, reflecting differences in system type, utilization and operations, and included life-cycle phases.

    \begin{figure}[htbp!]
        \centering
        \includegraphics[width=1\textwidth, trim={2cm 1cm 2.5cm 1.5cm}, clip]{0_plots/figure_LCA.pdf}
        \caption[LCA components for BSSs]{
            \textbf{\ac{lca} components for \acp{bss}.} 
            This figure summarises the main stages and impact categories considered in an \ac{lca}. Based on Luo et al.~\cite{Luo2019}, adapted and elaborated by the author.
        }
        \label{fig:LCA_bss}
    \end{figure}

    Manufacturing includes raw material extraction and production of bicycles, docking stations, and (when applicable) batteries. Manufacturing is often the largest single contributor in docked systems, accounting for around 50\% of total life-cycle emissions~\cite{Calan2024, Sun2022}. For \acp{e-b}, hardware production represents 61\% of life-cycle emissions in docked systems and 52\% in dockless systems~\cite{Luo2019}, with batteries contributing 20--30\% of embodied \acp{ghg}~\cite{Zhu2023}. Material choices also shape embodied impacts: for private bicycles, aluminum frames have been reported to carry higher production impacts than alternatives such as steel~\cite{Dave2010, Leunberger2010, Luna2016} or bamboo~\cite{Agyekum2017}.

    While cycling itself is practically emission-free, the use phase includes supporting operations, with rebalancing consistently emerging as a dominant contributor. Luo et al.~\cite{Luo2019} estimated that rebalancing accounts for 36\% of total life-cycle emissions in docked systems and up to 73\% in dockless fleets~\cite{Luo2019}. Impacts vary with fleet size, system scale, vehicle type, and redistribution frequency~\cite{Bonilla-Alicea2020}, with reported service-vehicle intensities ranging from 0.012~km of van travel per bike-km in docked systems~\cite{Felipe-Falgas2022} to 0.10--0.26~km in dockless fleets~\cite{Sun2022, Cazzola2020}. Vehicle technologies further shape emissions, from conventional vans~\cite{Chen2020LifeCycle, Hollingsworth2019} to electric vans or cargo bikes for battery collection and charging~\cite{Felipe-Falgas2022, Calan2024}. For electrically assisted systems, charging adds additional impacts: Sun et al.~\cite{Sun2022} reported electricity demands with climate impacts ranging from about 10--20~g~\ce{CO2}-eq/km in low-carbon grids to more than seven times higher in coal-dominated regions~\cite{Zhu2023, Calan2024}.  Maintenance is often underrepresented due to limited operator records~\cite{Calan2024}, but when included it can contribute around 10--15\% of total life-cycle \acp{ghg} in docked systems, with larger shares in dockless fleets~\cite{Sun2022, Chen2024, deBortoli2021Environmental, Luo2019}. In here, assumptions about component lifetimes are central, and improving durability and repairability can mitigate environmental impacts~\cite{Zhu2023, Calan2024}. Most \acp{lca} also incorporate modal substitution by assigning each trip a counterfactual mode and estimating avoided life-cycle emissions; outcomes depend critically on the share of car and ride-hailing trips displaced versus low-carbon modes~\cite{Krauss2022}. For instance, in Barcelona, Felipe-Falgas et al.~\cite{Felipe-Falgas2022} show that shared \acp{e-b} can even increase \acp{ghg} when they mainly replace \ac{pt} or walking rather than car trips.

    End-of-life impacts are often underrepresented due to data limitations~\cite{Calan2024} and typically represent only a small share of total life-cycle emissions~\cite{Krauss2022}. Beyond formal disposal, fleet attrition and oversupply can exacerbate impacts: Calan et al.~\cite{Calan2024} noted that losing around 20\% of bikes annually can increase total fleet emissions by as much as 25\%. A dramatic demonstration of this challenge was seen in Chinese cities during the 2017--2018 dockless bike boom, when ``bike-sharing graveyards'' (Figure~\ref{fig:bss_graveyard}) emerged as excess and abandoned bikes accumulated in massive lots~\cite{Sun2022, AmusingPlanet2018}. In these situations, low utilization prevented manufacturing emissions from being distributed across enough trips, and inadequate recycling meant much of the material value was lost~\cite{Luo2019, Sun2022}.

\begin{figure}    
    \centering
    \includegraphics[width=0.8\textwidth]{0_plots/graveyard.jpg}
    \caption[BSS graveyard in Xiamen, China]{
        \textbf{\ac{bss} graveyard in Xiamen, China.} 
        Photograph sourced from \textit{Amusing Planet}~\cite{AmusingPlanet2018} illustrating large-scale accumulation of shared bikes.
    }
    \label{fig:bss_graveyard}
\end{figure}


Overall, the environmental performance of \acp{bss} depends on system design, operational practices, and local mobility patterns. Benefits from avoided transport emissions materialize primarily when car substitution is substantial and operational burdens are minimized, particularly rebalancing.








